\documentclass[a4paper,twoside]{article}


% --- RUIMTE EN LAYOUT ---
\usepackage{setspace}

\usepackage{geometry} %size regulating
\geometry{
 a4paper,
 left=25mm,
 right=25mm,
 top=20mm,
 bottom=20mm,
 }

\usepackage{parskip} % for better spacing between paragraphs
\usepackage{flushend} % Alleen nodig als je laatste pagina van artikel wilt vullen met tekst, zodat de kolommen even lang zijn. Anders kan je dit beter weglaten.

% --- WISKUNDE ENSYMBOLEN ---
\usepackage{amsmath, amssymb, amsfonts, amsthm} % For nicer math format
\usepackage{mathrsfs} % For curly arrr
\usepackage{bm} % For bold symbols eg \bm{\nabla}
\usepackage{esint} % For bolletje in integrals
\usepackage{dsfont}
\usepackage{siunitx} % Voeg deze expliciet toe voor die graden en eenheden!
\usepackage{braket} % For bra-ket notation
\usepackage{mathtools}
\usepackage{url}
\allowdisplaybreaks

% --- OVERIGE ---

\usepackage{graphicx} % Required for inserting images
\usepackage[hidelinks]{hyperref} % geen kaders om linken
\usepackage{subcaption}
%\usepackage[dutch]{babel}
\usepackage[T1]{fontenc}
\usepackage[utf8]{inputenc}
\usepackage{lmodern}
\usepackage{multicol}


% --- Tikz packages ---
\usepackage{tikz}
\usetikzlibrary{angles,quotes, babel} % for pic (angle labels)
\usetikzlibrary{arrows.meta,decorations.markings,calc, decorations.pathmorphing, decorations.pathreplacing, positioning, shapes, shadings, patterns}
\usepackage{xcolor}
\usepackage{tikz-3dplot}
\usepackage[siunitx]{circuitikz}

\usepackage{etoolbox} % ifthen
\usepackage[outline]{contour} % glow around text
\usepackage{xfp} % higher precision (16 digits?)
\contourlength{1.1pt}



% --- THEOREM STIJLEN MET EXTRA RUIMTE ---

\theoremstyle{definition}
\newtheorem{postulate}{Postulaat}
\newtheorem{definition}{Definitie}
\newtheorem{derivation}{Afleiding}
\newtheorem{proposition}{Propositie}
\newtheorem{example}{Voorbeeld}

\theoremstyle{plain}
\newtheorem{theorem}{Theorem}


% --- EXTRA COMMANDO'S ---

\newcommand{\dbar}{{\mkern+3mu\mathchar'26\mkern-12mu d}}

% --- Titel en Introductie ---

\title{Complexe Analyse: Uitwerking test 1}
\author{T. Cools}
\date{2026-2027}

\begin{document}


\begin{abstract}
    Uitwerking van test 1 uit het jaar 2023-2024.
\end{abstract}

\section{Opgave 1}
\subsection{Mijn Oplossing}
\subsubsection{Paramatrisatie}
\begin{align*}
    \Gamma_1: \qquad & z_1(t) = \cos(t) -2 - i \sin(t) \qquad & 0 \leq t \leq 2\pi \\
    \Gamma_2: \qquad & z_2(t) = t-1 \qquad & 0 \leq t \leq 2 \\
    \Gamma_3: \qquad & z_3(t) = 2-\cos(t) - i \sin(t) \qquad & 0 \leq t \leq 2\pi \\
\end{align*}

\subsubsection{Totale lengte}
\[l(\Gamma) = 4\pi +2\]
\subsubsection{Rescaling + lijmen}

\[
z(t) = \begin{cases}
    \cos((4\pi +2)t)-2-i\sin((4\pi + 2)t) & \left(0 \leq t \leq \frac{2\pi}{4\pi + 2}\right) \\
    (4\pi + 2)\left(t-\frac{2\pi}{4\pi + 2}\right)-1 & \left(\frac{2\pi}{4\pi + 2} \leq t \leq \frac{2\pi + 2}{4\pi + 2}\right) \\
    2-\cos\left((4\pi +2)\left(t- \frac{2\pi + 2}{4\pi + 2}\right)\right)-i\sin\left((4\pi + 2)\left(t- \frac{2\pi + 2}{4\pi + 2}\right)\right) & \left(\frac{2\pi + 2}{4\pi + 2} \leq t \leq 1\right)
\end{cases}
\]
Vereenvoudigd:
\[
z(t) = \begin{cases}
    \cos((4\pi +2)t)-2-i\sin((4\pi + 2)t) & \left(0 \leq t \leq \frac{\pi}{\pi + 1}\right) \\
    (4\pi + 2)\left(t-\frac{\pi}{2\pi + 1}\right)-1 & \left(\frac{\pi}{2\pi + 1} \leq t \leq \frac{\pi + 1}{2\pi + 1}\right) \\
    2-\cos\left((4\pi +2)\left(t- \frac{\pi + 1}{2\pi + 1}\right)\right)-i\sin\left((4\pi + 2)\left(t- \frac{\pi + 1}{2\pi + 1}\right)\right) & \left(\frac{\pi + 1}{2\pi + 1} \leq t \leq 1\right)
\end{cases}
\]

\section{Opgave 2}
\subsection{Mijn oplossing}
\subsubsection{Stap 1: Partial fraction decomposition}

\[f(z) = \frac{2z^2 - z + 1}{(z-1)^2(z+2)} = \frac{A_0^{(0)}}{z+2} + \frac{A_0^{(1)}}{(z-1)^2} + \frac{A_1^{(1)}}{z-1}\]
Uitwerken\dots
\[f(z) = \frac{2z^2 - z + 1}{(z-1)^2(z+2)} = \frac{7}{z+2} + \frac{2}{3(z-1)^2} + \frac{11}{9(z-1)}\]
\subsubsection{Stap 2: Integraal}
\[
\int_\Gamma f(z)dz = \int_\Gamma\frac{7}{z+2}dz  + \int_\Gamma\frac{2}{3(z-1)^2} dz + \int_\Gamma\frac{11}{9(z-1) }dz
\]

We gebruiken eq. (4.52) uit de cursus:
\[
\int_{C_r} (z-z_0)^n dz = \begin{cases}
    0 \quad & n \neq -1\\
    2 \pi i \quad & n = -1
\end{cases}
\]
\begin{align*}
    \int_\Gamma f(z)dz &= \int_\Gamma\frac{7}{z+2}dz  + \int_\Gamma\frac{11}{9(z-1) }dz\\
    &= 7 \int_\Gamma\frac{1}{z+2}dz  +  \frac{11}{9}\int_\Gamma\frac{1}{z-1} dz\\
    &= 7 \cdot 2\pi i + (-1) \frac{11}{9} 2\pi i\\
    &= \pi i \frac{104}{9}\\
\end{align*}
\textbf{Let op: we hebben maal een factor -1 genomen doordat het punt 1 op de tekening in wijzerzin ligt. (zie figuur!)}
\end{document}
